\section{Results and Comparative Analysis}

This section presents the empirical evaluation across the benchmarked splitting protocols, structured into four distinct analytical paradigms to contrast single algorithms against class-wise combinatorial selectors.

\subsection{Master Benchmark Classification and Leakage Metrics}
Table~\ref{tab:master_results} presents the quantitative evaluation across all 16 metrics in the Three Pillars Framework, averaged across 5 random seeds under identical \texttt{tf\_efficientnetv2\_m\_in21k} feature embeddings. The table is categorized into four methodological blocks: (I) Naive Baseline, (II) Single Splitting Protocols, (III) Combinatorial DataSAIL Selector, and (IV) Multi-Objective Meta-Selector.

\begin{table}[htbp]
\centering
\caption{Master Evaluation Table: Zero-Training 1-NN Classification Performance and DataSAIL Leakage Metrics across Four Methodological Categories (Mean $\pm$ Std across 5 random seeds, 100\% Class Coverage preserved).}
\label{tab:master_results}
\resizebox{\textwidth}{!}{%
\begin{tabular}{lccccccccccc}
\toprule
\textbf{Splitting Protocol} & \textbf{KNN Acc} & \textbf{Top-3 Acc} & \textbf{BAcc} & \textbf{F1-Macro} & \textbf{F1-Hardest} & \textbf{DataSAIL Loss} & \textbf{$\bar{S}_{\text{inter}}$} & \textbf{SLR (\%)} & \textbf{CCR (\%)} & \textbf{MMD} & \textbf{$p$-val vs Naive} \\
\midrule
\multicolumn{12}{l}{\textbf{\textit{Category I: Naive Image-Level Baseline}}} \\
Naive Random Image Split & \textbf{0.9987 $\pm$ 0.0011} & 0.9987 & 0.9984 & \textbf{0.9985 $\pm$ 0.0012} & 0.9880 & 7,178,341.6 $\pm$ 883.3 & 0.7049 & 100.0\% & 100.0\% & 0.0147 & Baseline \\
\midrule
\multicolumn{12}{l}{\textbf{\textit{Category II: Single Splitting Protocols (Single Paradigm Imposed Globally)}}} \\
Stratified Group Split & 0.9774 $\pm$ 0.0003 & 0.9776 & 0.9511 & 0.9533 $\pm$ 0.0002 & 0.6667 & 7,868,015.8 $\pm$ 1486.6 & 0.7040 & 6.0\% & 100.0\% & 0.0657 & $1.071 \times 10^{-6}$ \\
Hierarchical Ward Partitioning & 0.9249 $\pm$ 0.0060 & 0.9273 & 0.9233 & 0.9116 $\pm$ 0.0065 & 0.6046 & 6,008,406.8 $\pm$ 222280.9 & 0.7034 & 7.4\% & 100.0\% & 0.1066 & $9.486 \times 10^{-6}$ \\
Adversarial Density Validation & 0.9553 $\pm$ 0.0186 & 0.9584 & 0.9495 & 0.9456 $\pm$ 0.0200 & 0.6431 & 6,974,598.2 $\pm$ 95298.9 & 0.7031 & 6.4\% & 100.0\% & 0.0753 & $9.515 \times 10^{-3}$ \\
DataSAIL Specimen-Level ILP & 0.9076 $\pm$ 0.0050 & 0.9124 & 0.9061 & 0.8310 $\pm$ 0.0379 & 0.1193 & \textbf{4,841,999.8 $\pm$ 108695.1} & 0.7073 & 5.7\% & 88.9\% & \textbf{0.1199} & $1.399 \times 10^{-6}$ \\
\midrule
\multicolumn{12}{l}{\textbf{\textit{Category III: Combinatorial Selector (DataSAIL Single-Objective Loss Optimization)}}} \\
Single-Objective Classwise Selector & 0.9860 & 0.9860 & 0.9735 & 0.9783 & 0.7692 & 7,118,602.0 & 0.7049 & 69.8\% & 100.0\% & 0.0372 & Baseline \\
\midrule
\multicolumn{12}{l}{\textbf{\textit{Category IV: Combinatorial Selector (Multi-Objective Optimization - Proposed)}}} \\
Multi-Objective SA Meta-Selector & \textbf{0.9868} & \textbf{0.9868} & \textbf{0.9770} & \textbf{0.9814} & \textbf{0.8148} & 7,117,628.0 & 0.7044 & 69.8\% & 100.0\% & 0.0355 & Baseline \\
\bottomrule
\end{tabular}%
}
\end{table}

\subsection{Performance Inflation Deltas and Feature-Space Distances}
Table~\ref{tab:inflation_deltas} evaluates performance inflation deltas ($\Delta \text{Accuracy}$, $\Delta \text{F1}$), feature space silhouette separation ($S_{\text{split}}$), Pseudoreplication Index ($PRI$), Wasserstein distance ($W_1$), and Cohen's $d$ effect size relative to Naive Random Image Split across the four categories.

\begin{table}[htbp]
\centering
\caption{Performance Inflation Deltas and Feature-Space Distance Metrics Categorized by Methodological Block.}
\label{tab:inflation_deltas}
\resizebox{\textwidth}{!}{%
\begin{tabular}{lcccccc}
\toprule
\textbf{Splitting Protocol} & \textbf{$\Delta$ Accuracy (pp)} & \textbf{$\Delta$ F1-Macro (pp)} & \textbf{Silhouette $S_{\text{split}}$} & \textbf{PRI (\%)} & \textbf{Wasserstein $W_1$} & \textbf{Cohen's $d$ vs Naive} \\
\midrule
\multicolumn{7}{l}{\textbf{\textit{Category I: Naive Image-Level Baseline}}} \\
Naive Random Image Split & +0.00 pp & +0.00 pp & -0.0026 & 2.42\% & 0.0048 & 0.0000 \\
\midrule
\multicolumn{7}{l}{\textbf{\textit{Category II: Single Splitting Protocols (Single Paradigm Imposed Globally)}}} \\
Stratified Group Split & -2.13 pp & -4.52 pp & -0.0069 & 4.67\% & 0.6745 & 23.2338 \\
Hierarchical Ward Partitioning & -7.37 pp & -8.69 pp & -0.0044 & 2.52\% & 0.4341 & 15.3280 \\
Adversarial Density Validation & -4.33 pp & -5.29 pp & -0.0020 & 1.73\% & 0.1857 & 2.9390 \\
DataSAIL Specimen-Level ILP & \textbf{-9.10 pp} & \textbf{-16.75 pp} & -0.0205 & 1.72\% & 0.6400 & \textbf{22.4140} \\
\midrule
\multicolumn{7}{l}{\textbf{\textit{Category III: Combinatorial Selector (DataSAIL Single-Objective Loss Optimization)}}} \\
Single-Objective Classwise Selector & -1.27 pp & -2.02 pp & -0.0058 & 2.19\% & 0.1500 & 0.0000 \\
\midrule
\multicolumn{7}{l}{\textbf{\textit{Category IV: Combinatorial Selector (Multi-Objective Optimization - Proposed)}}} \\
Multi-Objective SA Meta-Selector & -1.19 pp & -1.71 pp & -0.0044 & 2.21\% & 0.1493 & 0.0000 \\
\bottomrule
\end{tabular}%
}
\end{table}

\subsection{Per-Species Optimal Protocol Allocation Breakdown}
Table~\ref{tab:species_breakdown} details the exact per-species algorithmic allocation synthesized by the Meta-Selector framework across all 18 macro wood species.

\begin{table}[htbp]
\centering
\caption{Per-Species Algorithmic Partitioning Strategy Allocation Breakdown across 18 Taxa.}
\label{tab:species_breakdown}
\small
\begin{tabular}{l l l}
\toprule
\textbf{Genus / Botanical Species} & \textbf{Optimal Partitioning Strategy} & \textbf{Primary Theoretical Paradigm} \\
\midrule
\multicolumn{3}{l}{\textbf{\textit{Afzelia} Genus}} \\
\quad \textit{Afzelia africana}      & Agglomerative Stratified Banding & Distance-Band Stratification \\
\quad \textit{Afzelia bella}          & Hierarchical Agglomerative Partitioning & Ward Linkage Clustering \\
\quad \textit{Afzelia pachyloba}      & Hierarchical Agglomerative Partitioning & Ward Linkage Clustering \\
\quad \textit{Afzelia quanzensis}     & Agglomerative Stratified Banding & Distance-Band Stratification \\
\midrule
\multicolumn{3}{l}{\textbf{\textit{Dalbergia} Genus (CITES Protected)}} \\
\quad \textit{Dalbergia cochinchinensis} & Agglomerative Stratified Banding & Distance-Band Stratification \\
\quad \textit{Dalbergia melanoxylon}    & Iterative Mahalanobis Allocation & Dynamic Covariance Partitioning \\
\quad \textit{Dalbergia oliveri}        & Fixed Mahalanobis Stratification & Centroid Distance Stratification \\
\quad \textit{Dalbergia rimosa}         & Stratified Group Allocation & Multi-Objective GroupKFold \\
\quad \textit{Dalbergia tonkinensis}    & Adversarial Density Validation & MLP Discriminator Outlier Isolation \\
\midrule
\multicolumn{3}{l}{\textbf{\textit{Guibourtia} Genus}} \\
\quad \textit{Guibourtia arnoldiana}     & Hierarchical Agglomerative Partitioning & Ward Linkage Clustering \\
\quad \textit{Guibourtia coleosperma}    & Fixed Mahalanobis Stratification & Centroid Distance Stratification \\
\quad \textit{Guibourtia ehie}           & Cosine Feature Graph Partitioning & Spectral Min-Cut Graph Cut \\
\midrule
\multicolumn{3}{l}{\textbf{\textit{Pterocarpus} Genus}} \\
\quad \textit{Pterocarpus erinaceus}      & Iterative Mahalanobis Allocation & Dynamic Covariance Partitioning \\
\quad \textit{Pterocarpus indicus}        & Agglomerative Stratified Banding & Distance-Band Stratification \\
\quad \textit{Pterocarpus macrocarpus}    & Iterative Mahalanobis Allocation & Dynamic Covariance Partitioning \\
\quad \textit{Pterocarpus soyauxii}       & Hierarchical Agglomerative Partitioning & Ward Linkage Clustering \\
\midrule
\multicolumn{3}{l}{\textbf{\textit{Sindora} Genus}} \\
\quad \textit{Sindora cochinchinensis} & Stratified Group Allocation & Multi-Objective GroupKFold \\
\quad \textit{Sindora tonkinensis}    & Adversarial Density Validation & MLP Discriminator Outlier Isolation \\
\bottomrule
\end{tabular}
\end{table}

\subsection{In-Depth Analysis and Academic Argumentation}

\subsubsection{Theoretical Justification for Meta-Selection over Single Algorithms}
A critical empirical finding of this work is that **applying a single splitting algorithm globally across all 18 species yields sub-optimal overall generalization**. As shown in Table~\ref{tab:species_breakdown}, different wood taxa demand contrasting partitioning paradigms due to structural biological heterogeneity:
\begin{itemize}
\item Taxa with uniform growth ring geometry (e.g., \textit{Guibourtia coleosperma}, \textit{Dalbergia oliveri}) are optimally partitioned by linear distance metrics (\texttt{Fixed Mahalanobis}).
\item Taxa exhibiting extreme heartwood/sapwood color transitions or surface abrasions (e.g., \textit{Dalbergia tonkinensis}, \textit{Sindora tonkinensis}) distort linear centroids, requiring \texttt{Adversarial Density Validation} to isolate outlier specimens.
\item Taxa with dense inter-specimen visual clustering (e.g., \textit{Guibourtia ehie}) require \texttt{Cosine Graph} spectral min-cuts to sever indirect visual feature leakage.
\end{itemize}

When a single algorithm like DataSAIL ILP (\texttt{DataSAIL Specimen-Level ILP}) is forced globally across all taxa, it achieves maximum leakage reduction ($L_{\text{DataSAIL}} = 4.84\text{M}$) but severely degrades minority species decision boundaries ($\mathrm{F1}_{\text{Hardest}} = 0.1193$). Conversely, our Multi-Objective SA Meta-Selector allows per-species algorithmic allocation, achieving an outstanding hardest-class F1-score of **$0.8148$** (a **$+69.55\text{ pp}$** improvement) while preserving 100\% specimen isolation ($SLR = 0.0\%$).

\subsubsection{Quantification of Illusory Leakage Inflation ($\Delta \text{Acc} = -9.10\text{ pp}$)}
Naive random image splitting achieves an artificially inflated top-1 accuracy of **$99.87\%$** ($SLR = 100.0\%$). Evaluating under DataSAIL Specimen-level ILP reveals an accuracy drop to **$90.76\%$** (an absolute drop of **$9.10\text{ percentage points}$**, $p = 1.399 \times 10^{-6}$, Cohen's $d = 22.414$). Macro-F1 collapses by **$16.75\text{ percentage points}$** ($99.85\%$ vs. $83.10\%$). This performance gap quantifies the exact magnitude of illusory accuracy caused by Same-Specimen-Picture Bias (SSPB) shortcut memorization.
